\documentclass[11pt,a4paper]{article}

\usepackage[T1]{fontenc}
\usepackage[utf8]{inputenc}
\usepackage[english]{babel}
\usepackage{lmodern}
\usepackage{amsmath,amssymb}
\usepackage[margin=2.5cm]{geometry}
\usepackage{xcolor}
\usepackage{listings}
\usepackage{longtable,booktabs,array,calc}
\usepackage{enumitem}
\usepackage{hyperref}
\usepackage{newunicodechar}

% Unicode symbols used in the text (keeps the file pdfLaTeX-compatible)
\newunicodechar{°}{\ensuremath{^\circ}}
\newunicodechar{±}{\ensuremath{\pm}}
\newunicodechar{²}{\ensuremath{^2}}
\newunicodechar{·}{\ensuremath{\cdot}}
\newunicodechar{×}{\ensuremath{\times}}
\newunicodechar{θ}{\ensuremath{\theta}}
\newunicodechar{φ}{\ensuremath{\varphi}}
\newunicodechar{₀}{\ensuremath{_0}}
\newunicodechar{→}{\ensuremath{\rightarrow}}
\newunicodechar{∂}{\ensuremath{\partial}}
\newunicodechar{∇}{\ensuremath{\nabla}}
\newunicodechar{−}{\ensuremath{-}}
\newunicodechar{≈}{\ensuremath{\approx}}
\newunicodechar{≤}{\ensuremath{\leq}}
\newunicodechar{≥}{\ensuremath{\geq}}

\definecolor{codebg}{RGB}{246,247,249}
\definecolor{codekw}{RGB}{0,82,155}
\definecolor{codecm}{RGB}{110,120,110}
\definecolor{codestr}{RGB}{163,21,21}
\lstset{
  language=Python,
  basicstyle=\ttfamily\footnotesize,
  keywordstyle=\color{codekw}\bfseries,
  commentstyle=\color{codecm}\itshape,
  stringstyle=\color{codestr},
  backgroundcolor=\color{codebg},
  frame=single, rulecolor=\color{black!15},
  breaklines=true, postbreak=\mbox{$\hookrightarrow$\space},
  columns=fullflexible, keepspaces=true, showstringspaces=false,
  xleftmargin=2pt, xrightmargin=2pt,
  literate={°}{{$^\circ$}}1 {²}{{$^2$}}1 {−}{{-}}1 {×}{{$\times$}}1
           {≤}{{$\leq$}}1 {≥}{{$\geq$}}1 {→}{{$\rightarrow$}}1 {·}{{$\cdot$}}1
           {…}{{\ldots}}1 {–}{{--}}1 {—}{{---}}1 {“}{{``}}1 {”}{{''}}1
}
\lstset{literate={°}{{$^\circ$}}1 {²}{{$^2$}}1 {−}{{-}}1 {×}{{$\times$}}1
           {≤}{{$\leq$}}1 {≥}{{$\geq$}}1 {→}{{$\rightarrow$}}1 {·}{{$\cdot$}}1
           {…}{{\ldots}}1 {–}{{--}}1 {—}{{---}}1 {“}{{``}}1 {”}{{''}}1}

\providecommand{\tightlist}{\setlength{\itemsep}{0pt}\setlength{\parskip}{0pt}}
\setlength{\parindent}{0pt}
\setlength{\emergencystretch}{3em}
\sloppy
\setlength{\parskip}{4pt}
\hypersetup{colorlinks=true, linkcolor=codekw, urlcolor=codekw}

\title{L.A.N.D.\ Topography\\[2pt]\large Documentation}
\author{NASA Space Apps - Patos team - Luca Naranza}
\date{19 November 2026}

\begin{document}

\maketitle
\tableofcontents
\newpage

\hypertarget{overview}{%
\section{Overview}\label{overview}}

The topography pipeline turns LOLA/PGDA polar DEMs into three things per
candidate landing site: a 360° horizon profile, terrain statistics
(slope, relief, coverage) and a grid of 100 × 100 m cells that the
webapp heatmap consumes. Everything lives in
\texttt{Hackaton\_Topography/scripts/}.

\textbf{Data flow}

\begin{enumerate}
\def\labelenumi{\arabic{enumi}.}
\tightlist
\item
  \texttt{download\_dem.py} fetches the large regional
  mosaic (LDEM 60S, 240 m/px) with resumable HTTP Range requests.
\item
  \texttt{dem\_products.py} makes sure an elevation
  raster is used even when someone points the code at a slope product
  (\texttt{LDSM\_*}).
\item
  \texttt{site\_coordinates.py} validates the site
  catalog (lat/lon per site) and runs a preflight before any
  computation.
\item
  \texttt{analyze\_all.py} is the batch driver: for
  each of the 27 site folders in
  \texttt{dati\_siti\_topografia/} it launches
  \texttt{script.py} as a subprocess, then
  independently re-verifies the results and writes the summary tables.
\item
  \texttt{script.py} is the per-site worker: it calls
  \texttt{horizon.py} for the horizon and computes
  slope/aspect/cell statistics itself.
\item
  \texttt{finalize\_results.py} re-checks the whole
  output tree and writes \texttt{verifica\_finale.json}
  and the comparison chart.
\end{enumerate}

\textbf{Inputs per site} (5 m/px, polar stereographic, lunar sphere R =
1 737 400 m):

\begin{longtable}[]{@{}
  >{\raggedright\arraybackslash}p{(\columnwidth - 4\tabcolsep) * \real{0.3333}}
  >{\raggedright\arraybackslash}p{(\columnwidth - 4\tabcolsep) * \real{0.3333}}
  >{\raggedright\arraybackslash}p{(\columnwidth - 4\tabcolsep) * \real{0.3333}}@{}}
\toprule\noalign{}
\begin{minipage}[b]{\linewidth}\raggedright
Product
\end{minipage} & \begin{minipage}[b]{\linewidth}\raggedright
Meaning
\end{minipage} & \begin{minipage}[b]{\linewidth}\raggedright
How it's used
\end{minipage} \\
\midrule\noalign{}
\endhead
\bottomrule\noalign{}
\endlastfoot
\texttt{*\_surf.tif} & Elevation (m) & Main input for
horizon and slope \\
\texttt{*\_slp.tif} & Slope (°) from PGDA & Only as a
cross-check against the recomputed slope \\
\texttt{*\_toterr.tif} & Vertical uncertainty (m) &
Sanity range check \\
\texttt{*\_ldec.tif} & Number of LOLA shots per pixel &
0 means interpolated terrain, not missing data \\
\end{longtable}

\textbf{Outputs per site}
(\texttt{risultati topografia/Sito\_*}):
\texttt{horizon.json},
\texttt{horizon\_local.json},
\texttt{horizon.png},
\texttt{summary.json},
\texttt{site\_cells.csv},
\texttt{validation.json},
\texttt{README.md}, \texttt{run.log}.
Global: \texttt{riepilogo.json/csv},
\texttt{confronto*.png},
\texttt{verifica\_finale.json}.

\textbf{Example (Sito\_NPA, from its \texttt{run.log}
and \texttt{summary.json}):} 4000 × 4000 px grid (20 ×
20 km), no missing elevation, elevation −2612 to −1347 m, mean slope
6.47°, 80.05 \% of the area ≤ 10°, horizon between 0.69° and 4.72°.

Conventions used everywhere: azimuth 0° = geographic north, 90° = east;
longitudes east-positive in {[}−180, 180); body-fixed frame
MOON\_ME\_DE421; no atmospheric refraction (correct for the Moon).

\hypertarget{getting-the-data-download_dem.py-and-dem_products.py}{%
\section{\texorpdfstring{Getting the data: \texttt{download\_dem.py}
and
\texttt{dem\_products.py}}{Getting the data: download\_dem.py and dem\_products.py}}\label{getting-the-data-download_dem.py-and-dem_products.py}}

\hypertarget{download_dem.py-62-lines}{%
\subsection{\texorpdfstring{\texttt{download\_dem.py} (62
lines)}{download\_dem.py (62 lines)}}\label{download_dem.py-62-lines}}

A script with no functions at module level, run as
\texttt{python download\_dem.py <url> <destination>}.
It exists because the regional mosaic is several GB and a plain download
that dies at 90 \% starts over.

\begin{lstlisting}[language=Python]
11  url, destination = sys.argv[1:3]
12  path=Path(destination).resolve();path.parent.mkdir(parents=True,exist_ok=True)
13  part=path.with_suffix(path.suffix+'.part');checkpoint=path.with_suffix(path.suffix+'.download.json')
\end{lstlisting}

\begin{itemize}
\tightlist
\item
  \textbf{11--13} Read the two CLI arguments, create the destination
  folder, and define two sidecar files:
  \texttt{X.tif.part} (the file being written) and
  \texttt{X.tif.download.json} (which blocks are done).
\end{itemize}

\begin{lstlisting}[language=Python]
14  with urllib.request.urlopen(urllib.request.Request(url,method='HEAD'),timeout=60) as response:
15      total=int(response.headers['Content-Length']);modified=response.headers.get('Last-Modified')
16  if path.exists():
17      if path.stat().st_size != total:raise ValueError('File finale di dimensione inattesa')
18      print('File gia completo:',path,flush=True);sys.exit(0)
19  block=16*1024*1024
\end{lstlisting}

\begin{itemize}
\tightlist
\item
  \textbf{14--15} A \texttt{HEAD} request gets the
  total size and \texttt{Last-Modified} without
  downloading anything.
\item
  \textbf{16--18} If the final file already exists, it is accepted only
  if its size matches exactly; otherwise the script stops instead of
  overwriting.
\item
  \textbf{19} Work unit: 16 MiB blocks.
\end{itemize}

\begin{lstlisting}[language=Python]
20  if checkpoint.exists():
21      state=json.loads(checkpoint.read_text());assert state['url']==url and state['total_bytes']==total
22      completed=set(state['completed_blocks'])
23  else:
24      prefix=part.stat().st_size if part.exists() else 0
25      assert prefix <= total
26      completed=set(range(prefix//block))
27      state={'url':url,'total_bytes':total,'last_modified':modified,'block_bytes':block}
28      if shutil.disk_usage(path.parent).free < total-prefix+256*1024*1024:raise OSError('Spazio disco insufficiente')
29      with part.open('ab') as target:target.truncate(total)
\end{lstlisting}

\begin{itemize}
\tightlist
\item
  \textbf{20--22} Resume: if a checkpoint exists, it must refer to the
  same URL and the same size, then the set of finished blocks is
  reloaded.
\item
  \textbf{23--26} First run (or a \texttt{.part} from
  an older sequential downloader): every full block already in the
  prefix counts as done.
\item
  \textbf{28} Refuses to start without the missing bytes plus a 256 MiB
  margin of free disk.
\item
  \textbf{29} \texttt{truncate(total)} preallocates the
  \texttt{.part} to its final size, so blocks can be
  written at any offset in any order.
\end{itemize}

\begin{lstlisting}[language=Python]
30  def persist():
31      state['completed_blocks']=sorted(completed)
32      temp=checkpoint.with_suffix(checkpoint.suffix+'.tmp')
33      temp.write_text(json.dumps(state,indent=2),encoding='utf-8');temp.replace(checkpoint)
\end{lstlisting}

\begin{itemize}
\tightlist
\item
  \textbf{30--33} Atomic checkpoint: write to
  \texttt{.tmp}, then
  \texttt{replace()}. A crash mid-write never leaves a
  half-written JSON.
\end{itemize}

\begin{lstlisting}[language=Python]
35  def fetch(index):
36      start=index*block;end=min(total,start+block)-1
37      for attempt in range(4):
38          try:
39              request=urllib.request.Request(url,headers={'Range':f'bytes={start}-{end}'})
40              with urllib.request.urlopen(request,timeout=120) as response:
41                  assert response.status==206 and response.headers.get('Content-Range')==f'bytes {start}-{end}/{total}'
42                  data=response.read(end-start+2)
43                  assert len(data)==end-start+1,'Blocco incompleto'
44              with part.open('r+b') as target:target.seek(start);target.write(data)
45              return index
46          except Exception:
47              if attempt==3:raise
48              time.sleep(attempt+1)
\end{lstlisting}

\begin{itemize}
\tightlist
\item
  \textbf{36} Byte range of block \texttt{index}; the
  last block is shorter.
\item
  \textbf{41} The key safety check: the server must answer
  \texttt{206 Partial Content} with exactly the
  requested range. A server that ignores \texttt{Range}
  and returns \texttt{200} with the whole file would
  otherwise corrupt the output silently.
\item
  \textbf{42--43} Reads one byte more than expected, so a too-long body
  is detected as well as a short one.
\item
  \textbf{44} Each thread opens its own handle and writes at the block's
  offset.
\item
  \textbf{37, 46--48} Up to 4 attempts with linear back-off (1, 2, 3 s);
  the 4th failure propagates.
\end{itemize}

\begin{lstlisting}[language=Python]
49  indices=[i for i in range((total+block-1)//block) if i not in completed]
52  with concurrent.futures.ThreadPoolExecutor(max_workers=8) as pool:
53      for future in concurrent.futures.as_completed([pool.submit(fetch,i) for i in indices]):
54          completed.add(future.result());persist()
55          done=sum(min(block,total-i*block) for i in completed)
57  assert len(completed)==(total+block-1)//block
58  part.replace(path)
59  digest=hashlib.file_digest(path.open('rb'),'sha256').hexdigest()
\end{lstlisting}

\begin{itemize}
\tightlist
\item
  \textbf{49} Missing blocks = all blocks (ceiling division) minus the
  done ones.
\item
  \textbf{52--56} 8 parallel downloads; the checkpoint is updated after
  \textbf{every} finished block (from the main thread only, so no race
  on the JSON). Progress is computed from the real size of each block.
\item
  \textbf{57--58} Only when every block is present does
  \texttt{.part} become the final name.
\item
  \textbf{59--61} A SHA-256 of the result is stored in the checkpoint.
  Note: it is \emph{recorded}, not \emph{compared} against a published
  checksum. \texttt{hashlib.file\_digest} needs Python
  ≥ 3.11.
\end{itemize}

\hypertarget{dem_products.py-35-lines}{%
\subsection{\texorpdfstring{\texttt{dem\_products.py} (35
lines)}{dem\_products.py (35 lines)}}\label{dem_products.py-35-lines}}

One function, \texttt{resolve\_dem\_path(value)}. PGDA
distributes a slope map \texttt{LDSM\_60S\_240MPP\_ADJ}
next to the elevation map
\texttt{LDEM\_60S\_240MPP\_ADJ}. The horizon needs
\textbf{elevation}: a slope field cannot be integrated back to unique
heights. This function makes the swap explicit and safe.

\begin{lstlisting}[language=Python]
13      path = Path(value)
14      if not path.stem.upper().startswith('LDSM_'):
15          return path
16      match = re.fullmatch(r'LDSM_(\d+[NS])_(\d+MPP)_ADJ', path.stem, re.IGNORECASE)
\end{lstlisting}

\begin{itemize}
\tightlist
\item
  \textbf{14--15} Anything that isn't an LDSM file is passed through
  untouched.
\item
  \textbf{16} Extracts region (\texttt{60S}) and
  resolution (\texttt{240MPP}).
  \texttt{fullmatch} rejects related files such as
  error maps.
\end{itemize}

\begin{lstlisting}[language=Python]
21      expected = f'LDEM_{match[1]}_{match[2]}_ADJ'.upper()
22      matches = sorted({f.resolve() for f in path.parent.rglob('*')
23          if f.is_file() and f.suffix.lower() in ('.tif', '.tiff') and f.stem.upper() == expected})
\end{lstlisting}

\begin{itemize}
\tightlist
\item
  \textbf{21--23} Builds the twin name and searches the LDSM's folder
  \textbf{and its subfolders} (case-insensitive,
  \texttt{.tif}/\texttt{.tiff}). In
  this project it finds
  \texttt{dati\_regionali/60S/LDEM\_60S\_240MPP\_ADJ.TIF},
  as \texttt{run.log} shows.
\item
  \textbf{24--27} Zero matches or more than one is an error, never a
  guess.
\end{itemize}

\begin{lstlisting}[language=Python]
29      with rasterio.open(path) as slope, rasterio.open(dem) as elevation:
30          same_crs = (slope.crs is not None and elevation.crs is not None
31                      and slope.crs.to_dict() == elevation.crs.to_dict())
32          if (slope.count != 1 or elevation.count != 1 or slope.shape != elevation.shape
33                  or not slope.transform.almost_equals(elevation.transform) or not same_crs):
\end{lstlisting}

\begin{itemize}
\tightlist
\item
  \textbf{29--34} Confirms the two files really are twins: one band
  each, same pixel shape, same affine transform (within float
  tolerance), same CRS. Only then is the LDEM path returned.
\end{itemize}

\hypertarget{the-horizon-engine-horizon.py-246-lines}{%
\section{\texorpdfstring{The horizon engine: \texttt{horizon.py} (246
lines)}{The horizon engine: horizon.py (246 lines)}}\label{the-horizon-engine-horizon.py-246-lines}}

This is the physics core of the topography side. Its output
(\texttt{horizon.json}) is the contract the ephemeris
engine reads: for each integer azimuth 0--359°, the elevation angle
above which the Sun or Earth is visible from the site.

\hypertarget{constants-1415}{%
\subsection{Constants (14--15)}\label{constants-1415}}

\begin{lstlisting}[language=Python]
14  RADIUS_M = 1_737_400.0
15  LUNAR_CRS = CRS.from_proj4(f"+proj=longlat +R={RADIUS_M} +no_defs")
\end{lstlisting}

\begin{itemize}
\tightlist
\item
  \textbf{14} IAU mean lunar radius, the same sphere the PGDA products
  are defined on.
\item
  \textbf{15} A geographic (lon/lat) CRS on that sphere. Every lon/lat
  in the pipeline is expressed here; using EPSG:4326 (Earth WGS84) would
  silently put the points in the wrong place.
\end{itemize}

\hypertarget{demsampler-reading-heights-at-arbitrary-lonlat-1882}{%
\subsection{\texorpdfstring{\texttt{DEMSampler}: reading heights at
arbitrary lon/lat
(18--82)}{DEMSampler: reading heights at arbitrary lon/lat (18--82)}}\label{demsampler-reading-heights-at-arbitrary-lonlat-1882}}

\begin{lstlisting}[language=Python]
21      product_name = Path(source.name).stem.upper()
22      unit = (source.units[0] or '').lower() if source.count else ''
23      if product_name.startswith('LDSM_') or product_name.endswith('_SLP') or 'degree' in unit:
24          raise ValueError(...)
25      if source.count != 1 or source.crs is None:
28      if not (np.isclose(crs.ellipsoid.semi_major_metre, RADIUS_M)
29              and np.isclose(crs.ellipsoid.semi_minor_metre, RADIUS_M)):
31      self.to_xy = Transformer.from_crs(LUNAR_CRS, crs, always_xy=True)
32      self.data = self._read() if source.width * source.height <= 20_000_000 else None
\end{lstlisting}

\begin{itemize}
\tightlist
\item
  \textbf{21--24} Second line of defence against slope maps (after
  \texttt{dem\_products.py}): file names
  \texttt{LDSM\_*}, \texttt{*\_SLP},
  or a band unit containing ``degree'' are refused.
\item
  \textbf{25--30} Exactly one band, a CRS present, and a spherical body
  of radius 1 737 400 m (semi-major = semi-minor).
\item
  \textbf{31} Transformer lon/lat → projected x/y of this raster.
  \texttt{always\_xy=True} fixes the argument order to
  (lon, lat).
\item
  \textbf{32} Rasters up to 20 Mpx (the 4000×4000 site DEMs = 16 Mpx,
  128 MB as float64) are loaded whole. Bigger ones (the regional mosaic)
  use the tile cache below.
\end{itemize}

\begin{lstlisting}[language=Python]
36  def _read(self, window=None):
38      data = src.read(1, window=window, masked=True).astype("float64").filled(np.nan)
39      data = data * src.scales[0] + src.offsets[0]
40      data[~np.isfinite(data)] = np.nan
\end{lstlisting}

\begin{itemize}
\tightlist
\item
  \textbf{38} Reads band 1 with NoData masked, converts to float64,
  turns masked pixels into NaN.
\item
  \textbf{39} Applies the GeoTIFF scale/offset (raw integers → metres
  when the file stores them that way).
\item
  \textbf{40} Any inf becomes NaN: from here on, NaN is the single
  meaning of ``no data''.
\end{itemize}

\begin{lstlisting}[language=Python]
44      x, y = self.to_xy.transform(lon, lat)
45      inverse = ~self.source.transform
46      col = inverse.a*np.asarray(x) + inverse.b*np.asarray(y) + inverse.c
47      row = inverse.d*np.asarray(x) + inverse.e*np.asarray(y) + inverse.f
49      row, col = np.atleast_1d(row) - 0.5, np.atleast_1d(col) - 0.5
51      inside = (np.isfinite(row) & np.isfinite(col) & (row >= 0) & (col >= 0)
52                & (row <= self.source.height-1) & (col <= self.source.width-1))
55          result[indices] = map_coordinates(self.data, [row[indices], col[indices]], order=1, ...)
\end{lstlisting}

\begin{itemize}
\tightlist
\item
  \textbf{44} lon/lat → projected metres.
\item
  \textbf{45--47} The inverse affine transform turns metres into
  fractional pixel coordinates, written out explicitly so it works on
  whole arrays at once.
\item
  \textbf{49} The subtle line. Rasterio's (0,0) is the \emph{top-left
  corner} of the first pixel;
  \texttt{scipy.ndimage.map\_coordinates} treats
  integer indices as pixel \emph{centres}. Subtracting 0.5 aligns them.
  Without it every sample would be shifted by half a pixel (2.5 m
  locally, 120 m on the 240 m mosaic).
\item
  \textbf{51--52} Only points whose 2×2 bilinear neighbourhood is inside
  the raster are sampled; the rest stay NaN. Points in the outer
  half-pixel are therefore treated as outside.
\item
  \textbf{55} \texttt{order=1} = bilinear
  interpolation, \texttt{prefilter=False} because no
  spline is used. A NaN neighbour makes the result NaN, so no-data never
  leaks in as a number.
\end{itemize}

\begin{lstlisting}[language=Python]
61          tile_h, tile_w = self.source.block_shapes[0]
62          if tile_h == 1:
63              tile_h = 256
65          keys = (row[indices].astype(int)//tile_h)*ncols + col[indices].astype(int)//tile_w
66          for key in np.unique(keys):
71                  height = min(tile_h+1, self.source.height-r0)
72                  width = min(tile_w+1, self.source.width-c0)
76                  while self._cache_bytes > 64*1024*1024 and len(self._tiles)>1:
77                      _, removed = self._tiles.popitem(last=False)
79              self._tiles.move_to_end(cache_key)
\end{lstlisting}

\begin{itemize}
\tightlist
\item
  \textbf{61--63} Uses the file's internal tiling. Striped TIFFs (1-row
  blocks) get 256-row bands instead.
\item
  \textbf{65--66} Each sample is assigned to its tile; samples are
  processed tile by tile. The comment on 58--60 explains why: a ray 200
  km long crosses the raster diagonally, and one bounding window around
  it would read millions of useless pixels.
\item
  \textbf{71--72} Tiles are read with \textbf{one extra row and column},
  so the bilinear neighbour of a sample on the last row/column of the
  tile is present.
\item
  \textbf{76--79} LRU cache capped at 64 MiB: the oldest tile is evicted
  first, and a tile that is used again moves to the end. Rays from
  neighbouring azimuths reuse the same tiles near the site, so most
  reads hit the cache.
\end{itemize}

\hypertarget{elevation_angles-8589}{%
\subsection{\texorpdfstring{\texttt{elevation\_angles}
(85--89)}{elevation\_angles (85--89)}}\label{elevation_angles-8589}}

\begin{lstlisting}[language=Python]
87      d = np.asarray(distances_m, dtype=float)
88      return np.degrees(np.arctan2(
89          np.asarray(heights) - ground_m - observer_height_m - d*d/(2*RADIUS_M), d))
\end{lstlisting}

For a terrain point at surface distance \emph{d} with height \emph{h},
the apparent elevation seen from an eye at
\texttt{ground + observer\_height} is

\begin{quote}
e = atan( (h − h₀ − h\_obs − d²/2R) / d )
\end{quote}

The term \textbf{d²/2R} is how far the sphere curves away below the
local horizontal plane: 288 m at 10 km, 11.5 km at 200 km. This is why
distant terrain must be very high to matter. It is the small-angle
version of the exact spherical geometry. I checked the leftover error:
it is of order h·θ²/2 (about 13 m at 200 km for h ≈ −2 km), or about
0.004°, which is negligible. \texttt{arctan2} keeps the
sign right for points below the horizontal.

\hypertarget{calculate_horizon-92206}{%
\subsection{\texorpdfstring{\texttt{calculate\_horizon}
(92--206)}{calculate\_horizon (92--206)}}\label{calculate_horizon-92206}}

\textbf{Validation (102--121)}

\begin{itemize}
\tightlist
\item
  \textbf{102--103} lat/lon are mandatory. The DEM centre is
  deliberately \emph{not} used as a fallback (see Issues: this is why
  the current catalog blocks the run).
\item
  \textbf{104--110} Steps and distances must be positive and finite,
  with step ≤ near distance ≤ max distance; observer height ≥ 0.
\item
  \textbf{111--112} Redundant: lines 102--103 already reject a missing
  lat or lon.
\item
  \textbf{113--117} Remembers what was \emph{requested}, then maps every
  LDSM path to its LDEM.
\item
  \textbf{120--121} \texttt{describe\_site} (from
  \texttt{site\_coordinates.py}) checks that the point
  lies inside the near DEM and normalises lat/lon.
\end{itemize}

\textbf{Opening the sources (122--136)}

\begin{itemize}
\tightlist
\item
  \textbf{122--127} \texttt{ExitStack} keeps every
  raster open for the whole loop and closes them all even on error.
  Source order is: near DEM (index 0), far DEM (1), fallbacks (2\ldots).
\item
  \textbf{129--136} Writes a provenance record for each source (file,
  shape, resolution, scale/offset, bounds) into the metadata.
\end{itemize}

\textbf{Ground height and sample distances (137--145)}

\begin{lstlisting}[language=Python]
139      ground = float(near.sample([lon], [lat])[0])
142      distances = np.append(np.arange(step_m, near_distance_m, step_m), near_distance_m)
144          distances = np.concatenate((distances,
145              np.arange(near_distance_m+far_step_m, max_distance_m, far_step_m), [max_distance_m]))
\end{lstlisting}

\begin{itemize}
\tightlist
\item
  \textbf{137--138} Excludes the exact poles, where ``north'' is
  undefined.
\item
  \textbf{139} Ground height under the observer, sampled bilinearly from
  the 5 m DEM.
\item
  \textbf{142} Near samples: 5, 10, \ldots, 9995, 10000 m (2000
  samples).
\item
  \textbf{144--145} Far samples: 10 100 \ldots{} 200 000 m every 100 m
  (1900 samples). In total there are 3900 samples per azimuth × 360 ≈
  1.4 M height lookups per site.
\end{itemize}

\textbf{The ray loop (146--181)}

\begin{lstlisting}[language=Python]
146      geod = Geod(a=RADIUS_M, b=RADIUS_M)
150          lons, lats, _ = geod.fwd(np.full(distances.size, lon), np.full(distances.size, lat),
151                                    np.full(distances.size, az), distances)
\end{lstlisting}

\begin{itemize}
\tightlist
\item
  \textbf{146, 150--151} For azimuth \texttt{az},
  \texttt{Geod.fwd} walks along the \emph{great circle}
  starting at the site with true geographic bearing
  \texttt{az}, and returns lon/lat at every distance.
  This is what makes azimuth 0 = geographic north, not grid north. That
  matters at −87°, where the stereographic grid axes are rotated by up
  to tens of degrees relative to the meridians.
\end{itemize}

\begin{lstlisting}[language=Python]
154          local = distances <= near_distance_m
155          heights[local] = near.sample(lons[local], lats[local])
156          source_index[np.isfinite(heights)] = 0
157          local_angles = elevation_angles(heights, ground, observer_height_m, distances)
158          local_horizon.append(float(np.nanmax(local_angles)) ...)
\end{lstlisting}

\begin{itemize}
\tightlist
\item
  \textbf{154--156} First pass: only the 5 m DEM, only out to 10 km.
  Each sample is tagged with source 0.
\item
  \textbf{157--158} The \textbf{local-only horizon} is computed now,
  before far terrain is added. It is saved separately
  (\texttt{orizzonte\_locale}) so you can see how much
  the 200 km terrain adds.
\end{itemize}

\begin{lstlisting}[language=Python]
159          for index, sampler in enumerate(regional, 1):
160              missing = ~np.isfinite(heights)
163              samples = sampler.sample(lons[missing], lats[missing])
165              heights[positions] = samples
166              source_index[positions[np.isfinite(samples)]] = index
\end{lstlisting}

\begin{itemize}
\tightlist
\item
  \textbf{159--166} Each regional source fills \textbf{only} samples
  that are still NaN, in priority order. The 240 m mosaic therefore
  never overwrites 5 m data, and everything beyond 10 km comes from it.
  Gaps that no source covers stay NaN: they are never set to 0 m, which
  would invent a flat plain.
\end{itemize}

\begin{lstlisting}[language=Python]
170          angles = elevation_angles(heights, ground, observer_height_m, distances)
171          winner = int(np.nanargmax(angles))
172          horizons.append(float(angles[winner]))
\end{lstlisting}

\begin{itemize}
\tightlist
\item
  \textbf{170--172} The horizon at this azimuth is the \textbf{maximum
  elevation angle over all samples}: the highest terrain silhouette
  along the ray. \texttt{winner} remembers which sample
  it was.
\item
  \textbf{173--179} Coverage record per azimuth: valid/total samples,
  count per source, furthest valid distance, and whether coverage is
  complete.
\item
  \textbf{180--181} Obstruction record: distance, height and source of
  the point that sets the horizon. This is how you can tell whether a
  site's horizon is set by a nearby crater rim or a massif 150 km away.
\end{itemize}

\textbf{Result (182--206)}

\begin{itemize}
\tightlist
\item
  \textbf{182--189} Coverage counts as complete only if a far DEM exists
  \textbf{and} every azimuth has every sample. Warnings: no far DEM,
  incomplete coverage (horizon possibly underestimated), fallback used.
\item
  \textbf{190--206} The JSON contract: \texttt{lat},
  \texttt{lon},
  \texttt{orizzonte[360]},
  \texttt{orizzonte\_locale[360]}, and the metadata.
  The metadata include frame MOON\_ME\_DE421, the azimuth convention,
  the curvature model, and the explicit caveat that complete sample
  coverage does not guarantee every feature between samples was caught.
\end{itemize}

\hypertarget{plot_horizon-208241}{%
\subsection{\texorpdfstring{\texttt{plot\_horizon}
(208--241)}{plot\_horizon (208--241)}}\label{plot_horizon-208241}}

\begin{itemize}
\tightlist
\item
  \textbf{215--219} Reads the DEM decimated to about 1000 px on its
  longest side, averaging pixels
  (\texttt{Resampling.average}) so the overview image
  doesn't alias.
\item
  \textbf{220--226} Left panel: elevation map with the site (red +) and
  the DEM centre (grey ×). When the two coincide, as in the current
  results, that is visible here.
\item
  \textbf{231--233} Right panel: horizon versus azimuth, the near-only
  profile dashed and the full profile solid. Each array is extended by
  its first value (\texttt{[:1]}) so the curve closes
  at 360°.
\end{itemize}

\hypertarget{the-per-site-worker-script.py-277-lines}{%
\section{\texorpdfstring{The per-site worker: \texttt{script.py} (277
lines)}{The per-site worker: script.py (277 lines)}}\label{the-per-site-worker-script.py-277-lines}}

One script with two jobs, chosen by \texttt{--mode}:
\texttt{horizon} (calls
\texttt{horizon.py}), \texttt{metrics}
(slope, cells, summary) or \texttt{all}.
\texttt{analyze\_all.py} always runs it with
\texttt{--mode all --no-rasters}.

\hypertarget{statisticsvalues-2333}{%
\subsection{\texorpdfstring{\texttt{statistics(values)}
(23--33)}{statistics(values) (23--33)}}\label{statisticsvalues-2333}}

\begin{itemize}
\tightlist
\item
  \textbf{24} Drops NaN/inf first, so every statistic describes only
  valid pixels.
\item
  \textbf{25--26} An empty array returns \texttt{None},
  not a fake 0.
\item
  \textbf{27--33} count, min, max, mean, std (accumulated in float64
  even when the input is float32, which avoids rounding drift over 16 M
  pixels), and the 5th/50th/95th percentiles.
\end{itemize}

\hypertarget{terrain_derivativeselevation-dx-dy-3655}{%
\subsection{\texorpdfstring{\texttt{terrain\_derivatives(elevation,\ dx,\ dy)}
(36--55)}{terrain\_derivatives(elevation, dx, dy) (36--55)}}\label{terrain_derivativeselevation-dx-dy-3655}}

\begin{lstlisting}[language=Python]
42      dz_south, dz_east = np.gradient(elevation, dy, dx)
43      dz_north = -dz_south
\end{lstlisting}

\begin{itemize}
\tightlist
\item
  \textbf{42} \texttt{np.gradient} returns one
  derivative per array axis. Axis 0 is the row index, which grows
  \textbf{downwards = southwards} on a north-up grid, so the first
  output is ∂z/∂(south); axis 1 grows eastwards. In the interior it uses
  central differences: (z{[}i+1{]} − z{[}i−1{]}) / 2·dy. The slope is
  therefore measured over a \textbf{10 m baseline} at 5 m/px.
\item
  \textbf{43} Flips the sign to get ∂z/∂(north).
\end{itemize}

\begin{lstlisting}[language=Python]
44      valid = np.isfinite(elevation)
46      stencil[1:-1, 1:-1] = (
47          valid[1:-1, 1:-1] & valid[:-2, 1:-1] & valid[2:, 1:-1]
48          & valid[1:-1, :-2] & valid[1:-1, 2:]
49      )
\end{lstlisting}

\begin{itemize}
\tightlist
\item
  \textbf{44--49} A pixel is trusted only if it and its four neighbours
  (N, S, E, W) all have data. Edge pixels are always excluded, because
  \texttt{np.gradient} falls back to one-sided
  differences there. This is why NPA has 15 984 004 slope values =
  3998², not 4000².
\end{itemize}

\begin{lstlisting}[language=Python]
50      magnitude = np.hypot(dz_east, dz_north)
51      slope = np.degrees(np.arctan(magnitude))
52      aspect = np.degrees(np.arctan2(-dz_east, -dz_north)) % 360
54      aspect[~stencil | (magnitude < 1e-7)] = np.nan
\end{lstlisting}

\begin{itemize}
\tightlist
\item
  \textbf{50--51} Slope = atan(\textbar∇z\textbar), the steepest-descent
  angle, in degrees.
\item
  \textbf{52} Aspect = compass direction of the \textbf{downhill} vector
  (−∇z): \texttt{atan2(east, north)} gives a bearing
  clockwise from north, and \texttt{\% 360} maps it
  into {[}0, 360). Measured from \textbf{grid} north, not geographic
  north (stated in the docstring and in
  \texttt{summary.json}).
\item
  \textbf{53--54} Untrusted pixels become NaN. Aspect is also NaN on
  flat ground (\textbar∇z\textbar{} \textless{} 1e-7), where the
  direction is meaningless.
\end{itemize}

\hypertarget{lunar_lonlat-5861}{%
\subsection{\texorpdfstring{\texttt{lunar\_lonlat}
(58--61)}{lunar\_lonlat (58--61)}}\label{lunar_lonlat-5861}}

\begin{itemize}
\tightlist
\item
  \textbf{60--61} Converts projected x/y to lon/lat on the lunar sphere.
  The comment explains why: rasterio would happily transform to
  EPSG:4326, which is Earth's ellipsoid.
\end{itemize}

\hypertarget{analyze...-64195}{%
\subsection{\texorpdfstring{\texttt{analyze(...)}
(64--195)}{analyze(...) (64--195)}}\label{analyze...-64195}}

\textbf{Setup and checks (65--91)}

\begin{itemize}
\tightlist
\item
  \textbf{65--68} Resolves LDSM→LDEM. If lat/lon were given,
  \texttt{describe\_site} returns the site's pixel and
  position, and the coordinate source is recorded.
\item
  \textbf{71--76} One band, projected CRS in metres, lunar sphere. Line
  75 tests the radius with
  \texttt{to\_dict().get("R") != 1737400}, an exact
  match on a proj4 key. \texttt{horizon.py} does the
  same check more robustly with \texttt{np.isclose} on
  the ellipsoid (see Issues).
\item
  \textbf{78--79} The affine must have no rotation terms
  (\texttt{b = d = 0}), positive x step and negative y
  step: north-up. The north/south derivative naming in
  \texttt{terrain\_derivatives} depends on this.
\item
  \textbf{81--84} Reads the elevation as float32 (64 MB for 16 Mpx),
  applies scale and offset, and turns NoData into NaN.
\item
  \textbf{85} Slope and aspect for the whole grid.
\item
  \textbf{89} Cell size in pixels:
  \texttt{round(100/5) = 20} → 20 × 20 px cells.
  \texttt{round} means a tile size that isn't a
  multiple of the pixel is snapped; the actual size is reported as
  \texttt{tile\_size\_actual\_m}.
\end{itemize}

\textbf{Optional rasters (94--116)}, skipped in batch runs by
\texttt{--no-rasters}

\begin{itemize}
\tightlist
\item
  \textbf{95--108} \texttt{terrain\_metrics.tif}: 3
  float32 bands (elevation, slope, grid aspect), deflate compression,
  floating-point predictor, with band names and units.
\item
  \textbf{110--116} \texttt{slope\_filter.tif}: uint8
  with 1 = slope ≤ threshold, 0 = steeper, 255 = unknown.
\end{itemize}

\textbf{The 100 m cell loop (118--148)}

\begin{lstlisting}[language=Python]
119      for row in range(0, src.height, tile_rows):
120          for col in range(0, src.width, tile_cols):
121              end_row, end_col = min(row + tile_rows, src.height), min(col + tile_cols, src.width)
124              zv, sv = z[np.isfinite(z)], s[np.isfinite(s)]
125              if not zv.size:
126                  continue
127              x, y = affine * ((col + end_col) / 2, (row + end_row) / 2)
\end{lstlisting}

\begin{itemize}
\tightlist
\item
  \textbf{119--121} Walks the grid in 20 × 20 px steps. Cells at the
  edge are clipped (none occur here, since 4000 / 20 = 200).
\item
  \textbf{125--126} Cells with no elevation at all are dropped from the
  CSV.
\item
  \textbf{127} Cell centre in projected metres. The pixel indices are
  corner-based, so (col + end\_col)/2 is exactly the geometric centre.
\end{itemize}

Fields written per cell (128--140):

\begin{longtable}[]{@{}
  >{\raggedright\arraybackslash}p{(\columnwidth - 4\tabcolsep) * \real{0.3333}}
  >{\raggedright\arraybackslash}p{(\columnwidth - 4\tabcolsep) * \real{0.3333}}
  >{\raggedright\arraybackslash}p{(\columnwidth - 4\tabcolsep) * \real{0.3333}}@{}}
\toprule\noalign{}
\begin{minipage}[b]{\linewidth}\raggedright
Field
\end{minipage} & \begin{minipage}[b]{\linewidth}\raggedright
Formula
\end{minipage} & \begin{minipage}[b]{\linewidth}\raggedright
Note
\end{minipage} \\
\midrule\noalign{}
\endhead
\bottomrule\noalign{}
\endlastfoot
\texttt{elevation\_coverage\_pct} & valid z / all
pixels & \\
\texttt{slope\_coverage\_pct} & valid slope / all
pixels & Cells on the DEM border are \textless{} 100 \% \\
\texttt{elevation\_mean\_m} & mean of valid z & \\
\texttt{relief\_m} & max − min
(\texttt{np.ptp}) & Includes any regional tilt; not
roughness \\
\texttt{slope\_mean\_deg},
\texttt{slope\_max\_deg} & over valid slope & \\
\texttt{passes\_slope\_filter\_pct} & count(s ≤ 10°) /
\textbf{all} pixels & NaN compares False, so missing data counts as
``fail'' \\
\end{longtable}

\begin{itemize}
\tightlist
\item
  \textbf{141--148} Converts all centres to lunar lon/lat in a single
  call and writes \texttt{site\_cells.csv}. This is the
  file the webapp heatmap loads.
\end{itemize}

\textbf{Site and summary (150--195)}

\begin{itemize}
\tightlist
\item
  \textbf{150--154} When a site is given: slope and aspect at the site's
  own pixel, plus the record of the 100 m cell that contains it. That
  cell is found by snapping the pixel's row/col down to a multiple of
  20.
\item
  \textbf{155--156} DEM centre in x/y and lon/lat. For NPA that is x =
  −88 500 m, y = 29 500 m → −86.924°, −71.565°.
\item
  \textbf{157} Fixed thresholds 5/10/15/20° plus the user's own,
  deduplicated and sorted.
\item
  \textbf{168--171} Areas are pixel counts × 25 m². These are
  \emph{projected} areas; see Issues for the stereographic scale.
\item
  \textbf{173--177} For each threshold: area at or below it, and
  percentage of valid-slope pixels. For NPA ≤ 10° that is 319.9 km² =
  80.05 \%. Unlike the per-cell \%, this denominator excludes NaN.
\item
  \textbf{180--183}
  \texttt{top\_cells\_by\_slope\_only}: only fully
  covered cells, sorted by pass \% (descending), then max slope, then
  relief (ascending), keeping the top 10. The name is honest: it is a
  slope ranking, not a landing score.
\item
  \textbf{184--192} The limitations list written into every summary.
\item
  \textbf{194} \texttt{allow\_nan=False}: if a NaN ever
  slipped into the summary, the write would fail instead of producing
  invalid JSON.
\end{itemize}

\hypertarget{main-198268}{%
\subsection{\texorpdfstring{\texttt{main()}
(198--268)}{main() (198--268)}}\label{main-198268}}

\begin{itemize}
\tightlist
\item
  \textbf{199--221} CLI. Defaults point to Site01 and
  \texttt{risultati topografia/Sito\_01}; every horizon
  parameter is exposed (\texttt{--observer-height},
  \texttt{--near-distance},
  \texttt{--step}, \texttt{--far-dem},
  \texttt{--max-distance},
  \texttt{--far-step},
  \texttt{--fallback-dem} repeatable).
\item
  \textbf{222--223} Horizon and all modes refuse to run without
  \texttt{--lat} and \texttt{--lon}.
\item
  \textbf{224--237} \texttt{resolved()} maps LDSM→LDEM
  for each DEM argument and records every substitution. That record
  produces the \texttt{LDSM -> LDEM:} line in
  \texttt{run.log} and the
  \texttt{prodotti\_risolti} metadata.
\item
  \textbf{238--244} Files must exist, 0 ≤ max slope \textless{} 90, tile
  size \textgreater{} 0.
\item
  \textbf{245--262} Horizon: compute, attach the product mapping and
  coordinate source, write \texttt{horizon.json}
  (strict JSON), plot, print min/max (NPA: 0.694--4.721°) and any
  warnings.
\item
  \textbf{263--267} Metrics: \texttt{analyze()}, then
  print the key part of the summary. That printout is what appears in
  \texttt{run.log}.
\end{itemize}

Note: the NPA \texttt{summary.json} currently on disk
has \textbf{no \texttt{site} key}. It was produced by
an earlier version that used the DEM centre, before the
explicit-coordinate logic was added (see Issues).

\hypertarget{orchestration-and-verification-site_coordinates.py-analyze_all.py-finalize_results.py}{%
\section[Orchestration and verification]{\texorpdfstring{Orchestration and verification:
\texttt{site\_coordinates.py}, \texttt{analyze\_all.py},
\texttt{finalize\_results.py}}{Orchestration and verification: site\_coordinates.py, analyze\_all.py, finalize\_results.py}}\label{orchestration-and-verification-site_coordinates.py-analyze_all.py-finalize_results.py}}

\hypertarget{site_coordinates.py-175-lines}{%
\subsection{\texorpdfstring{\texttt{site\_coordinates.py} (175
lines)}{site\_coordinates.py (175 lines)}}\label{site_coordinates.py-175-lines}}

This module keeps two things that are easy to confuse apart:
\textbf{where the site is} (the catalog) and \textbf{where the DEM tile
is centred}.

\textbf{\texttt{checked\_coordinates} (20--25)}

\begin{itemize}
\tightlist
\item
  \textbf{21--22} Rejects anything non-numeric, \emph{including
  booleans}: in Python \texttt{True} is an
  \texttt{int}, so \texttt{lat: true}
  in JSON would otherwise pass as 1°.
\item
  \textbf{23} lat strictly inside (−90, 90); lon accepted in {[}−180,
  360{]} so that both conventions work.
\item
  \textbf{25} \texttt{(lon+180)\%360-180} wraps to
  {[}−180, 180): 350° → −10°.
\end{itemize}

\textbf{\texttt{dem\_geometry} (28--42)}: same CRS
checks as \texttt{horizon.py}
(\texttt{np.isclose} on the ellipsoid). Then the centre
of the raster (\texttt{width/2, height/2} through the
affine transform) is converted to lunar lon/lat. Used for reporting
only.

\textbf{\texttt{describe\_site(dem, lat, lon)}
(45--77)}, called by \texttt{horizon.py},
\texttt{script.py} and
\texttt{analyze\_all.py}:

\begin{lstlisting}[language=Python]
51          x,y=to_xy.transform(lon,lat)
54          row,col=src.index(x,y)
58          c=inv.a*x+inv.b*y+inv.c-.5
59          r=inv.d*x+inv.e*y+inv.f-.5
60          if not (0<=r<=src.height-1 and 0<=c<=src.width-1) or min(src.shape)<2:
62          ri,ci=min(math.floor(r),src.height-2),min(math.floor(c),src.width-2)
63          wr,wc=r-ri,c-ci
64          patch=src.read(1,window=Window(ci,ri,2,2),masked=True)...
68          ground=float((1-wr)*(1-wc)*patch[0,0]+wr*(1-wc)*patch[1,0]+(1-wr)*wc*patch[0,1]+wr*wc*patch[1,1])
70      _,_,distance=Geod(a=RADIUS_M,b=RADIUS_M).inv(geometry['dem_center_lon'],geometry['dem_center_lat'],lon,lat)
\end{lstlisting}

\begin{itemize}
\tightlist
\item
  \textbf{51, 54--56} Projects to x/y, finds the containing pixel
  (\texttt{src.index} floors), and refuses points
  outside the DEM.
\item
  \textbf{58--61} Centre-based fractional pixel coordinates (the same
  −0.5 convention as \texttt{DEMSampler}). A site in
  the outer half-pixel ring is refused, because it has no full 2×2
  neighbourhood.
\item
  \textbf{62--68} Reads exactly the 2×2 patch and does bilinear
  interpolation \textbf{by hand}. Weight \texttt{wr}
  goes to the lower row (\texttt{patch[1,…]}) and
  \texttt{wc} to the right column. The
  \texttt{min(…, height-2)} clamp handles a point lying
  exactly on the last centre row/column. This value must agree with
  \texttt{DEMSampler.sample} to 1e-6 m (checked in
  \texttt{verify}), so two independent implementations
  cross-check each other.
\item
  \textbf{70} Geodesic distance from the DEM centre to the site: the
  number that shows whether the analysis point is the tile centre or a
  real chosen site.
\end{itemize}

\textbf{\texttt{load\_catalog} (80--94)}

\begin{itemize}
\tightlist
\item
  \textbf{82} \texttt{utf-8-sig} tolerates a BOM, which
  Windows Notepad adds.
\item
  \textbf{83--84} The catalog must declare frame MOON\_ME\_DE421, units
  degree, east-positive longitude. This makes the reference frame a
  deliberate statement rather than an assumption.
\item
  \textbf{88--93} Every entry needs a non-empty
  \texttt{site\_id}; IDs are matched case-insensitively
  and duplicates are an error.
\item
  \textbf{94} Returns the entries plus the SHA-256 of the catalog file.
  The hash is written into every output, so each result can be traced to
  the exact coordinate file that produced it.
\end{itemize}

\textbf{\texttt{find\_dem} (97--100)}: exactly one
\texttt{*\_surf.tif} per site folder.

\textbf{\texttt{create\_template} (103--111)}: writes a
catalog with every site folder and
\texttt{lat/lon/source = null}. It \textbf{never
overwrites} an existing file (105). This is the state the repository is
in now.

\textbf{\texttt{preflight} (114--160)}

\begin{itemize}
\tightlist
\item
  \textbf{118--139} Per folder: find the DEM, compute the DEM centre,
  look up the catalog entry, validate lat/lon, require a non-empty
  \texttt{source}, then run
  \texttt{describe\_site}. The status is
  \texttt{READY},
  \texttt{MISSING\_COORDINATES} (null values) or
  \texttt{INVALID\_COORDINATES\_OR\_DEM}.
\item
  \textbf{140} \texttt{ready} only if every site is
  READY: it is all or nothing.
\item
  \textbf{143--159} Writes
  \texttt{controllo\_coordinate/coordinate\_preflight.json},
  a \texttt{.csv} and a
  \texttt{README.md} table that puts site lat/lon next
  to DEM-centre lat/lon and the distance between them.
\item
  \textbf{163--172} CLI: exit code 0 if ready, 2 otherwise.
\end{itemize}

\hypertarget{analyze_all.py-374-lines-the-batch-driver}{%
\subsection{\texorpdfstring{\texttt{analyze\_all.py} (374 lines): the
batch
driver}{analyze\_all.py (374 lines): the batch driver}}\label{analyze_all.py-374-lines-the-batch-driver}}

\textbf{Globals (19--31)}

\begin{itemize}
\tightlist
\item
  \textbf{24} \texttt{FAR\_REFERENCE} =
  \texttt{dati\_regionali/LDSM\_60S\_240MPP\_ADJ.tiff}.
  The LDSM file name only identifies the product; it is resolved to the
  LDEM.
\item
  \textbf{25--31} \texttt{FAR\_DEM},
  \texttt{FALLBACK\_DEMS},
  \texttt{MAX\_DISTANCE},
  \texttt{SITE\_CATALOG} and
  \texttt{SITE\_CATALOG\_HASH} are set in
  \texttt{main()} and read by
  \texttt{verify()} and
  \texttt{index()}. The ``NOT USED'' heading on 29 is
  misleading: the catalog \emph{is} used.
\end{itemize}

\textbf{\texttt{sample\_bilinear(source, lons, lats)}
(38--59)}: a second bilinear sampler for the regional DEM, deliberately
\textbf{not} using scipy, so the verification doesn't share code with
the thing it verifies.

\begin{itemize}
\tightlist
\item
  \textbf{44--45} Same −0.5 centre convention.
\item
  \textbf{48} Strict \texttt{< height-1}, where
  \texttt{DEMSampler} uses
  \texttt{<=}. The two differ only for a sample exactly
  on the last centre row/column.
\item
  \textbf{51--54} Reads one window spanning every requested sample: the
  bounding box of the whole ray. On the 240 m mosaic that is at most
  about 830 × 830 px for a 200 km ray, which is acceptable.
\item
  \textbf{55--58} Explicit 4-weight bilinear.
\end{itemize}

\textbf{\texttt{verify(dem, output, expected\_site)}
(62--232)}: everything a reviewer would check by hand.

\begin{itemize}
\tightlist
\item
  \textbf{66--68} Reloads the full 5 m elevation as float64.
\item
  \textbf{69--75} Consistency between files:

  \begin{itemize}
  \tightlist
  \item
    grid shape matches \texttt{summary.json};
  \item
    \texttt{horizon.json} used this exact DEM;
  \item
    lat/lon in \texttt{horizon.json} equal the catalog
    position to 1e-12°;
  \item
    \texttt{summary.site} exists, with the same lat/lon
    and pixel;
  \item
    ground height agrees to 1e-6 m.
  \end{itemize}

  Line 73 is why the existing results, which have no
  \texttt{site} key, cannot pass
  \texttt{verify} today.
\item
  \textbf{78--79} Exactly 360 finite horizon values.
\item
  \textbf{81--84} \textbf{Independent slope}: every 4th pixel in each
  direction (≈ 1 M points), central differences written out by hand:
  \texttt{(z[c+1]−z[c−1])/2dx}.
\item
  \textbf{86--99} For each of \texttt{surf},
  \texttt{slp}, \texttt{toterr},
  \texttt{ldec}: one band, identical shape, transform
  and CRS parameters; statistics; percentage missing. A CRS that is
  numerically equal but carries a different label becomes a warning, not
  a failure (the Site23 case).
\item
  \textbf{100--108} Checks that \texttt{\_slp} lies in
  {[}0, 90{]}, then compares the recomputed slope with PGDA's own slope:
  mean and max absolute difference go into
  \texttt{validation.json}. This is the real external
  check on \texttt{terrain\_derivatives}.
\item
  \textbf{109--115} Uncertainty and LDEC must be ≥ 0. For LDEC it
  records the \% of zero pixels (interpolated, never measured) and warns
  about \textbf{fractional} LDEC values: a shot count should be an
  integer (the Haworth case).
\item
  \textbf{120--131} A third bilinear sampler, this one on the in-memory
  5 m array.
\item
  \textbf{133--134} Rebuilds the distance array from the metadata.
\item
  \textbf{135--138} Ground agrees to 1e-7 m; the far DEM and the max
  distance are the expected ones.
\item
  \textbf{139} Local horizon ≤ full horizon at every azimuth. Adding
  terrain can only raise the maximum, so a violation would mean a bug.
\item
  \textbf{140--142} Per-azimuth source counts add up to the valid
  samples, and the \texttt{copertura\_completa} flag is
  consistent.
\item
  \textbf{147--158} \textbf{Independent horizon} at 0, 45, \ldots, 315°:
  same geodesic rays, same fill order, independent samplers, and
  \texttt{arctan} instead of
  \texttt{arctan2} (equal for d \textgreater{} 0). It
  must agree with \texttt{horizon.json} to \textless{}
  1e-7°. This checks the \emph{implementation} (sampling, fill order,
  indexing), not the \emph{model}: both use the same d²/2R formula.
\item
  \textbf{161--168} CSV row count ≤ ceil(H/20)·ceil(W/20), and equal to
  it when the DEM has no NoData. For a 4000² DEM that is 40 000 cells.
\item
  \textbf{175--179} Any warning makes the status
  \texttt{PASS\_WITH\_NOTES}; it is then saved.
\item
  \textbf{180--198} Builds the summary row. Notable fields:

  \begin{itemize}
  \tightlist
  \item
    \texttt{azimut\_ostruzione\_aumentata}: how many
    directions far terrain raised by \textgreater{} 1e-4°;
  \item
    \texttt{aumento\_massimo\_orizzonte\_deg}: the
    largest such increase;
  \item
    \texttt{distanza\_sito\_centro\_dem\_m}.
  \end{itemize}
\item
  \textbf{199--227} The per-site \texttt{README.md}.
\end{itemize}

\textbf{\texttt{index(rows, failures)} (235--275)}:
writes \texttt{riepilogo.json} (parameters, including
the catalog hash, plus all rows and errors) and
\texttt{riepilogo.csv}. If Excel has that CSV open
(Windows file lock → \texttt{PermissionError}), it
writes \texttt{riepilogo\_aggiornato.csv} instead
(240--243). Lines 267--268 hard-code notes about Site23 and Haworth;
they are not derived from the data.

\textbf{\texttt{main()} (278--369)}

\begin{itemize}
\tightlist
\item
  \textbf{288--296} Preflight first.
  \texttt{--preflight-only} just reports. If any site
  is not READY, the run \textbf{stops before touching any output}. This
  is what happens today.
\item
  \textbf{300--307} Resolves LDSM → LDEM, checks that the files exist,
  and requires a max distance ≥ 10 km.
\item
  \textbf{312--319} For each site folder (sorted case-insensitively):
  find the single \texttt{\_surf} DEM and its catalog
  entry.
\item
  \textbf{320--326} \texttt{--resume}: an existing
  result is reused only if the far DEM, max distance, fallbacks,
  coordinate origin, catalog hash and lat/lon (to 1e-12) all match. Any
  change in the catalog forces a recompute.
\item
  \textbf{327--334} Otherwise it runs
  \texttt{script.py --mode all --no-rasters …} as a
  \textbf{separate process}. Each site's 16 Mpx arrays are freed when
  the process exits, which matters on a laptop processing 27 sites.
  stdout and stderr go to \texttt{run.log}; a non-zero
  exit becomes an error.
\item
  \textbf{335--341} Stamps the catalog hash into
  \texttt{horizon.json} and writes
  \texttt{horizon\_local.json} (the near-only profile
  as its own file).
\item
  \textbf{342--347} \texttt{verify}; any exception is
  recorded per site with the full traceback in
  \texttt{error.log}, and the loop continues with the
  next site.
\item
  \textbf{348} \texttt{index()} after \textbf{every}
  site, so a crash at site 20 still leaves a valid summary of the first
  19.
\item
  \textbf{349--367} Two comparison charts:
  \texttt{confronto.png} (mean slope, \% ≤ 10°) and
  \texttt{confronto\_orizzonti.png} (far-terrain
  contribution).
\end{itemize}

\hypertarget{finalize_results.py-88-lines}{%
\subsection{\texorpdfstring{\texttt{finalize\_results.py} (88
lines)}{finalize\_results.py (88 lines)}}\label{finalize_results.py-88-lines}}

A last audit of the entire output tree, run after the batch.

\begin{itemize}
\tightlist
\item
  \textbf{16--17} Refuses to run if
  \texttt{riepilogo.json} has no catalog hash, which
  means results made from DEM centres. These are never ``finalised''.
\item
  \textbf{18--20} No errors, and the set of result sites equals the set
  of input folders.
\item
  \textbf{24--36} \texttt{normalize} renames the legacy
  source label \texttt{mosaico87S} to
  \texttt{lontano} inside
  \texttt{horizon.json} and
  \texttt{validation.json}. It is a leftover from an
  earlier run that used the 87S mosaic (whose
  \texttt{download.json} is still in
  \texttt{dati\_regionali}).
\item
  \textbf{30--31} All 8 files exist and are non-empty for every site.
\item
  \textbf{39--51} Re-asserts the contract:

  \begin{itemize}
  \tightlist
  \item
    PASS status;
  \item
    explicit coordinates with a matching catalog hash;
  \item
    lat/lon equal the validated position;
  \item
    360 finite values, with full ≥ local horizon;
  \item
    complete coverage to exactly 200 km;
  \item
    far DEM is \texttt{LDEM\_60S\_240MPP\_ADJ.TIF};
  \item
    no LDSM file was ever used as a height source, and the LDSM→LDEM
    mapping was recorded;
  \item
    no \texttt{error.log}.
  \end{itemize}
\item
  \textbf{55--61} Restores the batch globals from the report and
  rewrites the index. \textbf{62--67} then rewrites every
  \texttt{riepilogo*.csv} with the same rows; this
  duplicates what \texttt{index()} just did.
\item
  \textbf{69--75} Redraws
  \texttt{confronto\_orizzonti.png}.
\item
  \textbf{76--83} \texttt{verifica\_finale.json}: site
  count, total cells, 360 × sites directions, the largest
  independent-check error, and every site that has notes.
\end{itemize}

\hypertarget{issues-found-in-the-topography-code}{%
\section{Issues found in the topography
code}\label{issues-found-in-the-topography-code}}

Ranked by how much they affect the results. Overall, the code is
careful: frames, NoData, half-pixel offsets, LDSM/LDEM and coverage are
all handled explicitly and cross-checked. The weak point is not the
maths but the \textbf{state of the inputs}.

\hypertarget{blocking}{%
\subsection{Blocking}\label{blocking}}

\textbf{1. Every site coordinate in the catalog is
\texttt{null}, so the current results are horizons from
the DEM tile centre.}

\begin{itemize}
\tightlist
\item
  \texttt{config/site\_coordinates.json} is the
  untouched template: \texttt{lat},
  \texttt{lon} and \texttt{source} are
  \texttt{null} for all 27 sites.
\item
  The results in \texttt{risultati topografia/} were
  produced by an earlier version that used the DEM centre
  (\texttt{AVVERTENZA\_COORDINATE.md} says so). The NPA
  \texttt{summary.json} has no
  \texttt{site} key.
\item
  With the current code, \texttt{analyze\_all.py} stops
  at preflight, \texttt{verify()} would fail on line
  73, and \texttt{finalize\_results.py} refuses to run.
\item
  In practice: the horizon profiles the ephemeris engine reads are for
  the tile centre, which is not necessarily the chosen landing point. At
  −87° the Sun stays within about ±1.5° of the horizon and NPA's horizon
  reaches 4.7°, so moving the observer a few km can flip illumination
  for hours.
\item
  \textbf{Fix:} fill the catalog with the agreed points and a source
  (for example the Artemis III candidate-region coordinates the team
  used), then run \texttt{analyze\_all.py} and
  \texttt{finalize\_results.py}.
\end{itemize}

\hypertarget{affects-interpretation}{%
\subsection{Affects interpretation}\label{affects-interpretation}}

\textbf{2. Aspect is relative to grid north, not geographic north.}

\begin{itemize}
\tightlist
\item
  In polar stereographic, grid north at a point is rotated from true
  north by the point's longitude offset from the projection's central
  meridian. At NPA (lon ≈ −71.6°) that is a rotation of tens of degrees.
\item
  The code documents it in the docstring, the band name and the
  limitations list. But anyone using \texttt{aspect} to
  mean ``this slope faces the Sun or Earth'' gets a wrong direction.
\item
  \textbf{Fix:} add the meridian convergence to the aspect (computable
  per pixel from \texttt{pyproj}
  \texttt{Proj.get\_factors().meridian\_convergence}),
  or simply don't publish aspect.
\end{itemize}

\textbf{3. Horizon azimuth sampling is 1°.}

\begin{itemize}
\tightlist
\item
  At 200 km one degree is a 3.5 km arc, and at 50 km it is 0.9 km. A
  narrow peak that falls between two rays is missed.
\item
  When the ephemeris code interpolates between azimuths, it effectively
  assumes the silhouette is smooth.
\item
  For sunrise and sunset timing at the pole this is probably the largest
  geometric error left. Consider 0.25--0.5°, at a 2--4× runtime cost.
\end{itemize}

\textbf{4. Cell \texttt{passes\_slope\_filter\_pct}
counts missing slope as ``fail''} (\texttt{script.py}
139), while the DEM-wide \texttt{pct\_of\_valid\_slope}
excludes it.

\begin{itemize}
\tightlist
\item
  Border cells therefore look worse than they are.
\item
  It is conservative and commented, but the two percentages are not
  directly comparable.
\end{itemize}

\textbf{5. \texttt{relief\_m} is max − min.}

\begin{itemize}
\tightlist
\item
  On a uniformly tilted 100 m cell at 10°, relief is about 18 m with
  zero roughness.
\item
  For hazard ranking, detrended RMS (residual after fitting a plane)
  separates roughness from tilt.
\end{itemize}

\hypertarget{minor-or-cosmetic}{%
\subsection{Minor or cosmetic}\label{minor-or-cosmetic}}

\begin{itemize}
\tightlist
\item
  \textbf{6. Projected areas.} Polar stereographic with true scale at
  the pole has k = 2/(1+sin\textbar φ\textbar) ≈ 1.0007 at 87°. Areas
  are overestimated by about 0.14 \% and slopes underestimated by about
  0.07 \%. Negligible, and the code says so.
\item
  \textbf{7. Radius check inconsistency.}
  \texttt{script.py} 75 tests
  \texttt{crs.to\_dict().get("R") != 1737400} (exact
  match on a proj4 key). The other modules use
  \texttt{np.isclose} on the ellipsoid. A valid lunar
  CRS written as WKT without an \texttt{R} key would be
  rejected only by \texttt{script.py}.
\item
  \textbf{8. Redundant checks and stale text.}

  \begin{itemize}
  \tightlist
  \item
    \texttt{horizon.py} 111--112 repeats the check of
    102--103.
  \item
    The ``NOT USED'' comment at
    \texttt{analyze\_all.py} 29 is wrong.
  \item
    \texttt{index()} hard-codes notes about Site23 and
    Haworth, and parameter values (step 5 m, 10 km, 10°) that are not
    read from the run.
  \item
    \texttt{finalize\_results.py} rewrites the CSVs
    right after \texttt{index()} has written them.
  \end{itemize}
\item
  \textbf{9. The ``independent'' horizon check verifies the
  implementation, not the model.} It reuses the same distances, the same
  geodesic and the same d²/2R formula. A wrong model (say, a sign error
  in the curvature term) would pass. A one-off test against an exact
  spherical formula would close this gap. My own estimate of the
  approximation error is about 0.004° at 200 km, so the model itself is
  fine.
\item
  \textbf{10. The 5 m step along rays samples, it does not search.} A
  sub-5 m spike between samples is missed. The code states this
  (\texttt{limite\_risoluzione}), and at 5 m/px the DEM
  has no such detail anyway.
\item
  \textbf{11. \texttt{download\_dem.py}:}

  \begin{itemize}
  \tightlist
  \item
    it stores \texttt{Last-Modified} but never compares
    it on resume, so a file that changed on the server mid-download
    would produce a mixed file;
  \item
    the SHA-256 is recorded but not checked against a published value;
  \item
    if the checkpoint is deleted after preallocation, the full-size
    \texttt{.part} is taken as complete (line 26).
  \end{itemize}
\item
  \textbf{12. LDEC = 0 is interpolated terrain, not NoData.} It is
  handled correctly, but the \% of LDEC-zero pixels should be shown next
  to any per-cell result. A smooth-looking cell may simply be
  unmeasured.
\end{itemize}

\hypertarget{verified-correct-things-that-often-go-wrong-and-dont-here}{%
\subsection{Verified correct (things that often go wrong and don't
here)}\label{verified-correct-things-that-often-go-wrong-and-dont-here}}

\begin{itemize}
\tightlist
\item
  Half-pixel offset between the rasterio and scipy conventions
  (\texttt{horizon.py} 49, mirrored in all three
  samplers).
\item
  \texttt{np.gradient} axis order and the south→north
  sign flip.
\item
  Lunar sphere everywhere; never EPSG:4326.
\item
  Geographic azimuth through \texttt{Geod.fwd}, not
  grid bearings.
\item
  Far DEM only fills gaps; it never overwrites 5 m data. Missing terrain
  stays NaN, never 0.
\item
  LDSM is never used as heights; the mapping is logged and asserted.
\item
  Strict JSON (\texttt{allow\_nan=False}), per-site
  subprocess isolation, \texttt{index()} after every
  site, and resume keyed to the catalog hash.
\end{itemize}

\end{document}
